% Generated draft: requires mathematical and source review.
\documentclass[11pt]{article}
\usepackage[margin=25mm]{geometry}
\usepackage{fontspec}
\setmainfont{Latin Modern Roman}
\usepackage{amsmath,amssymb,mathtools,unicode-math}
\setmathfont{Latin Modern Math}
\usepackage{enumitem,needspace,xurl,hyperref}
\setlength{\parindent}{0pt}
\setlength{\parskip}{5pt}
\setlength{\emergencystretch}{3em}
\newcommand{\N}{\mathbb{N}}
\newcommand{\Z}{\mathbb{Z}}
\newcommand{\Q}{\mathbb{Q}}
\newcommand{\R}{\mathbb{R}}
\newcommand{\C}{\mathbb{C}}
\newcommand{\F}{\mathbb{F}}
\newcommand{\A}{\mathbb{A}}
\newcommand{\PP}{\mathbb{P}}
\newcommand{\E}{\mathbb{E}}
\newcommand{\eps}{\varepsilon}
\newcommand{\dd}{\,\mathrm{d}}
\begin{document}
\section*{UnsolvedMath ID: 30000303}
\textbf{Source problem code:} catalogue entry
\subsection{Definitions and mathematical statement}
Let $G_{s,\Delta}$ be the number of gaps in the best local alignment in the sequence-based random model of the report. The model, alignment score, gap penalty, sequence-length regime, and meaning of the parameters $s$, $\Delta$, and $\lambda$ are not specified in the supplied extraction. The stated asymptotic regime is
\[
 s\longrightarrow\infty,\qquad\Delta\longrightarrow\infty,
 \qquad s e^{-\lambda\Delta}\longrightarrow\beta\in(0,\infty).
\]
The conjecture asks whether, under this regime and the intended model assumptions,
\[
 G_{s,\Delta}\xrightarrow{\mathrm d}\operatorname{Poisson}(\mu),
 \qquad 0<\mu<\infty.
\]
Here $\mu$ may depend on $\beta$ and the fixed model data; no formula is supplied. A Poisson law has
\[
 \Pr(G=k)=e^{-\mu}\frac{\mu^k}{k!},\qquad k\in\N_0.
\]
The usual joint-limit interpretation would require this convergence along every admissible parameter sequence satisfying the displayed scaling, with a common limiting mean for fixed limiting data. The record does not explicitly spell out that uniformity or the allowed dependence of $\lambda$ and $\mu$.

The restored question compares the sequence-based model with a heuristic Poisson model in which the source says this conclusion holds. No definition of that heuristic model or transfer argument is supplied. The original extraction instead referred to an unstated asymptotic E-value formula. The supplied verification expressly replaces that paraphrase with the gap-count conjecture; the two targets are not identified as equivalent here. The dated literature assessment continues to discuss the E-value assertion and therefore requires reconciliation with the restored question.

To obtain a precise random variable one must also specify whether gaps count runs, inserted or deleted letters, or another unit, and how a best alignment is selected when maximizers are tied. No tie-breaking rule or distributional hypothesis on the input sequences is invented.
\subsection{Short English statement}
Under the stated coupled large-parameter scaling, the question asks whether the gap count in an optimal sequence-based local alignment converges to a Poisson law with a finite, strictly positive mean. The stochastic model and gap-counting conventions are missing, and older E-value metadata concern a different paraphrase.
\subsection{Sources}
\begin{itemize}
\item Local UnsolvedMath registry, permanent ID 30000303, code catalogue entry, ``Poisson Gap Counts in Gapped Local Alignments.'' Proposed year 2005; proposer unspecified.
\item \emph{Mathematical Population Genetics} (2006), Oberwolfach Reports, pp.\ 2241--2304, DOI 10.4171/owr/2005/40; original report \url{https://doi.org/10.4171/owr/2005/40}. The displayed report year 2006 differs from the DOI and proposed year 2005.
\item Supplied reference ``Approximate p-values for local sequence alignments'': \url{https://doi.org/10.1214/aos/1015951993}.
\item Statement status corrected\_verified. The supplied verification says the official report was checked directly and an extracted E-value paraphrase replaced by the actual conjecture on the asymptotic number of gaps. This source-check claim is not independently verified here.
\item Literature triage checked 2026-08-21: open. Its assessment says the gapped-alignment E-value assertion is conjectural or model-dependent outside restricted regimes and later work offers approximations and special cases rather than a general proof. Its applicability to the restored gap-count target requires checking. Created 2026-08-25; updated 2026-09-16.
\item Registry curation by Ulam AI / UnsolvedMath Contributors, CC BY 4.0; underlying report and literature provenance remain as supplied. Unreviewed statement-only draft.
\end{itemize}
\end{document}
